\chapter{Introdução}\label{cap:introducao}


%\chapter*[Introdução]{Introdução}
%\addcontentsline{toc}{chapter}{Introdução}

%Este documento segue as normas estabelecidas pela~\citeonline[3.1-3.2]{NBR6028:2003}.
A segurança pública hoje é um tema muito noticiado em veículos de mídia, segundo estudos hoje no Brasil aproximadamente 10\% dos crimes cometidos são solucionados \cite{um} e isso se dá principalmente devido ao fato das polícias terem um baixo efetivo não dando conta de cobrirem certas áreas. Outro fator muito importante é a falta de investimento em tecnologias o que não acontece em países de primeiro mundo que possuem uma alta taxa de soluções em casos criminais pelo fato destes países terem um alto investimento tecnológico na área de segurança. Um bom exemplo seria a China que possui hoje o melhor sistema de vigilância existente, composto de um complexo produto de IA utilizando visão computacional, o que faz com que eles possuam uma taxa de crimes mais baixa que a de países da Europa e equivalente ao de países como a Suíça \cite{dois}.  

No Brasil, os crimes são solucionados, pelo fato da polícia geralmente adquirir provas através de alguma imagem feita por uma câmera de terceiro, ou seja, equipamentos privados que foram instalados em residências ou empresas [16]. Contudo, percebemos que no Brasil existe uma grande falta de investimentos em pesquisa focada em novas tecnologias que poderiam suprir o mercado e principalmente a área de segurança. 

Hoje existem tecnologias com distribuição livre que podem ser utilizadas no mercado Brasileiro na área de segurança. Um exemplo são as ferramentas de visão computacional como o OpenCV, ferramenta muito poderosa que é utilizada em vários países com foco principalmente em reconhecimento facial, identificação de padrões de comportamento, detecção de objetos, entre outras várias funcionalidades, todas utilizando aprendizagem de máquina (Machine Learning) [17].  

Com investimento em pesquisa e desenvolvimento, a tecnologia de visão computacional pode ser aplicada como ferramenta para suprir a demanda de tecnologias na área de segurança pública no Brasil.

Alguns países que utilizam detecção facial já estão desenvolvendo sistemas integrados desta tecnologia com a utilização de um veículo aéreo não tripulado (VANT) \cite{dois} [2]. Com um VANT é possível cobrir uma grande área utilizando o espaço aéreo e em casos de desaparecimentos, este equipamento poderá acessar áreas de difícil acesso. Um bom exemplo e o caso do rompimento da barragem de Brumadinho aonde ocorreu de vários corpos não serem localizados pelo fato da lama dificultar as buscas, e uma empresa estrangeira veio com uma solução utilizando VANTs que possibilitaram a busca em certos locais, esses equipamentos possuíam câmeras com tecnologias que possibilitavam identificar um corpo em meio ou dentro da lama [18].  

%Tecnologias de drones para visão computacional ainda estão em desenvolvimento e possuem algumas limitações como a comunicação de uma controladora de vôo com um computador que possa sustentar a tarefa de processar todo o sistema de imagem e desempenhar comandos de vôo para o equipamento de vôo.

No site da Ardupilot \cite{tres}[3], já existe um projeto em desenvolvimento para conectar e configurar um computador complementar,  o Nvidia Jetson\footnote{\url{https://www.nvidia.com/pt-br/autonomous-machines/embedded-systems/jetson-nano/}}, Google Coral Dev Board\footnote{\url{https://coral.ai/products/dev-board}} e Raspberry Pi\footnote{\url{https://www.raspberrypi.org/products/raspberry-pi-3-model-b/}} a uma controladora de voo, essa conexão utiliza um protocolo chamado MAVLink e possui uma documentação bem desenvolvida que explica detalhadamente como controlar um VANT utilizando a ferramenta Dronekit, e a partir dele é possível prototipar uma grande diversidade de projetos unindo ferramentas e tecnologias diversas ao VANT.

Com isso surgiu a ideia do estudo de integrar via protocolo de comunicação (MAVLink) à Raspberry Pi com a controladora de voo do VANT, embarcar um sistema Linux com ferramentas de visão na Raspberry Pi, e assim possibilitar o desenvolvimento de um sistema com IA focado em visão computacional que possibilite o desenvolvimento de controles direcionais que atuem no VANT, sendo assim possível controla-lo através de um sistema que tenha uma câmera como sensor de orientação e direcionamento.  


% ---
\chapter{Objetivos e Justificativas}\label{cap:objJust}
% ---
A ideia surgiu devido ao interesse em integrar um veículo aéreo não tripulado, no caso um drone do tipo quadricóptero, com algum dispositivo para realização de alguma tarefa, e após pesquisar várias trabalhos, protótipos experimentais, bibliografia sobre o assunto e sites especializados em tecnologias aplicadas em Vants, como exemplo a visão computacional, se descobriu que é possível conectar através de um protocolo de comunicação, uma controladora de drone [7] a um computador complementar ou como é mais abordado mini computador [8], e enviar comandos para o Vant, criando um sistema de controle e automação de Vant através de visão computacional [3]. Em seguida também se descobriu que e possível integrar a ferramenta OpenCV na Raspberry Pi fazendo a integração de visão computacional com a tecnologia de um Vant.
Como o foco principal do trabalho é o controle utilizando visão computacional, primeiro se realizou uma pesquisa para descobrir se existe um computador complementar que pudesse comportar um software de visão computacional e que também tivesse capacidade computacional de compilar e rodar o algoritmo, e com processamento o suficiente para processar a tecnologia de visão computacional.
O Raspberry Pi foi a melhor opção devido a alguns fatores como; custo, tamanho, tempo.
A segunda pesquisa abordou como comunicar o computador complementar com a controladora do Vant e se existia a possibilidade de enviar comandos para a controladora de voo. 
Foi preciso buscar na internet muito material principalmente tutorias para aprender e dominar o funcionamento das ferramentas que seriam utilizadas no desenvolvimento, exemplo é o OpenCV, essa ferramenta e muito vasta possui muitas funcionalidades e módulos, então foi preciso descobrir qual seria a melhor maneira de utiliza-la no projeto, quais módulos utilizar, logo foi gasto muito tempo praticando e testando esses módulos para adequar quais seriam utilizados.  
Uma boa parte do trabalho foi realizando uma pesquisa bibliográfica sobre Deep Learning para o entendimento de como funcionam as tecnologias de visão computacional.

%-
\section{Objetivos Gerais}
%-
O objetivo geral deste trabalho é desenvolver o protótipo de um VANT controlado por um sistema embarcado com visão computacional para reconhecimento de padrões de comportamento de multidões, com reconhecimento facial, que possa ser empregado como uma ferramenta de segurança, podendo auxiliar em casos policiais, vigilância patrimonial, filmagens áreas, buscas em mata fechada, fiscalização de fronteira e aplicações civis nas quais seria inviável um ser humano trabalhar, realizando tarefas arriscadas e servindo como verdadeira ferramenta de trabalho.
Acredita-se que com o uso de imagens aéreas possa ser possível reconhecer e emitir alertas de possíveis suspeitos de crimes, e com isso buscar minimizar um pouco a falta de tecnologias que hoje são necessárias para ajudar no combate a violência. 
Um sistema de visão computacional integrado a um drone que através do desenvolvimento de um algoritmo e o treinamento de redes neurais, seja capaz de cumprir objetivos específicos, tendo como principal, o de identificar indivíduos através de reconhecimento facial, padrões de comportamento, e objetos específicos portados, e logo tomar a ação de rastrear o elemento que foi identificado através da movimentação aérea de um drone.   

%-
%%%%%%%%%%%%%%%%%%%%%%%%%%%%%%%%%%%%%%%%%%%%%%%%%%%%%%%%%%
%Normalmente não há problemas em usar caracteres acentuados em arquivos
%bibliográficos (\texttt{*.bib}). Porém, como as regras da ABNT fazem uso quase
%abusivo da conversão para letras maiúsculas, é preciso observar o modo como se
%escreve os nomes dos autores. Na~\autoref{tabela-acentos} você encontra alguns
%exemplos das conversões mais importantes. Preste atenção especial para `ç' e `í'
%que devem estar envoltos em chaves. A regra geral é sempre usar a acentuação
%neste modo quando houver conversão para letras maiúsculas.

%\begin{table}[htbp]
%\caption{Tabela de conversão de acentuação.}
%\label{tabela-acentos}

%\begin{center}
%\begin{tabular}{ll}\hline\hline
%acento & \textsf{bibtex}\\
%à á ã & \verb+\`a+ \verb+\'a+ \verb+\~a+\\
%í & \verb+{\'\i}+\\
%ç & \verb+{\c c}+\\
%\hline\hline
%\end{tabular}
%\end{center}
%\end{table}
%%%%%%%%%%%%%%%%%%%%%%%%%%%%%%%%%%%%%%%%%%%%%%%%%%%%%%%%%%%%%%%

\subsection{Objetivos Específicos}
%-
\begin{itemize}
    \item Pesquisar e estudar o funcionamento da ferramenta OpenCV;
    \item Criar e treinar uma rede neural;
    \item Treinar a ferramenta inserindo fotos para reconhecimento facial;
    \item Treinar a ferramenta para reconhecimento de padrões de comportamento;
    \item Comunicar com protocolo MAVLink a Raspberry Pi e a Pixhawk;
    \item Realizar testes de comunicação entre a Pixhawk e a Raspberry Pi;
    \item Desenvolver um sistema de tolerância a falhas de comunicação;
    \item Pesquisar e estudar como a Raspberry Pi envia comandos de vôo;
    \item Desenvolver um algoritmo que interprete os dados extraídas da visão computacional e os converta em comandos de voo;
    \item Desenvolver um algoritmo que envie comandos de voo para o drone;
\end{itemize}
%-

\subsection{Resultados Esperados}
%-
\begin{itemize}
    \item Perfeito funcionamento da ferramenta OpenCV;
    \item Que o algoritmo de rede neural seja compilado e executado pela Raspberry Pi;
    \item Reconhecimento das pessoas inseridas no algoritmo de reconhecimento facial;
    \item Distinção entre pessoas inseridas no algoritmo de rede neural, e as que não foram inseridas;
    \item Perfeita comunicação entre a Raspberry Pi e a Pixhawk;
    \item Perfeita tolerância de falhas do sistema;
    \item Desempenho satisfatório da Raspberry Pi em enviar comandos de vôo para o drone;
    \item Comportamento de voo do drone de maneira esperada que ele desempenhe;
    \item Perfeito funcionamento do algoritmo que será desenvolvido para o funcionamento do sistema que controla e corrige em que direção o drone deve seguir;
    \item Perfeito funcionamento dos sistemas de (software e o hardware) em seguir o objeto ou pessoa que foi inserido(a) no algoritmo de reconhecimento facial.
\end{itemize}
%-
%%%%%%%%%%%%%%%%%%%%%%%%%%%%%%%%%%%%%%%%%%%%%%%%%%%%%%%%%%%%%%%%%%%
%Consulte a FAQ com perguntas frequentes e comuns no portal do \abnTeX:
%\url{https://code.google.com/p/abntex2/wiki/FAQ}.

%Inscreva-se no grupo de usuários \LaTeX:
%\url{http://groups.google.com/group/latex-br}, tire suas dúvidas e ajude
%outros usuários.
%%%%%%%%%%%%%%%%%%%%%%%%%%%%%%%%%%%%%%%%%%%%%%%%%%%%%%%%%%%%%%%%%%%

\section{Justificativa} 
%-
Devido ao aumento da criminalidade e dos altos índices de violência no cotidiano das pessoas, surgiu a cultura do medo e o sentimento de insegurança. Tais fatores demandaram algumas mudanças nos serviços de segurança patrimonial bem como, nas formas de monitoramento. Nesse sentido, tornou-se necessário expandir as formas de controle, seja por meio de câmeras de vigilância ou monitoramento eletrônico [4]. No entanto, os equipamentos atuais como as câmeras utilizadas na segurança já não mais eficientes e possuem tecnologias ultrapassadas. 
Na cidade do Rio de janeiro existe um sistema de câmeras de alta tecnologia com capacidade para realizar reconhecimento facial. Este sistema foi implantado para testes através de uma parceria entre a OI e a Huawei [5], sendo que as câmeras e a tecnologia de reconhecimento facial embarcados são de propriedade da Chinesa Huawei e possuem alto custo monetário, e o que eu quero dizer com isso é que alguns estados Brasileiros passam por uma crise financeira, porem precisam investir em segurança [6].
Este contexto de insegurança pública e falta de investimento nacional em tecnologias avançadas de segurança é que originou a motivação para a proposta deste trabalho, ou seja, a necessidade da integração entre software e hardware de baixo custo, sendo aplicados para melhorar o desempenho do sistema de segurança no Brasil, visando a vigilância, através do sensoriamento utilizando câmeras e hardware de custo acessível e softwares livres.

%-





%\section*{Figuras}\label{sec:figuras}
%\addcontentsline{toc}{section}{figuras}

%As normas da~\citeonline[3.1-3.2]{NBR6028:2003} especificam que o caption da figura %deve vir abaixo da mesma.

%A Figura~\ref{fig:log} ilustra...

%\begin{figure}[htpb]
 %  \centering
 %  \includegraphics[scale=.3]{figs/logo}
 %  \caption{Breve explicação sobre a figura. Deve vir abaixo da mesma.}
 %  \label{fig:nome dado a figura}
%   \legend{Fonte: do autor}
%\end{figure}

%\section*{Tabelas}\label{sec:tabelas}
%\addcontentsline{toc}{section}{tabelas}

%A Tabela~\ref{tab:tabela} apresenta os resultados...

%\begin{table}[htpb]
   %\centering
   %\caption{Breve explicação sobre a tabela. Deve vir acima da %mesma.}\label{tab:tabela}
   %\begin{tabular}{|l|c|c|c|c|c|c|r|}
        %\hline
        %\small{XX} & \small{FF} & \small{PP} & \small{YY} & \small{Yr} & \small{xY} & %\small{Yx} & \small{ZZ} \\ \hline
               %615 &    18      &     2558   &    0,9930  &    0,9930  &    0,9930  & %   0,9930  &    0,9930  \\ \hline
               %615 &    18      &     2558   &    0,9930  &    0,9930  &    0,9930  & %   0,9930  &    0,9930  \\ \hline
               %615 &    18      &     2558   &    0,9930  &    0,9930  &    0,9930  & %   0,9930  &    0,9930  \\ \hline
               %615 &    18      &     2558   &    0,9930  &    0,9930  &    0,9930  &  %  0,9930  &    0,9930  \\ \hline
               %615 &    18      &     2558   &    0,9930  &    0,9930  &    0,9930  & %   0,9930  &    0,9930  \\ \hline
%   \end{tabular}
%\end{table}

%\section*{Motivação}\label{sec:motivacao}
%\addcontentsline{toc}{section}{Motivação}

%\lipsum[35]

